Las tres diferencias sucesivas de las primeras aproximaciones son \(0.03\), \(0.0075\) y \(0.001875\): se reducen por un factor próximo a cuatro. Bajo la hipótesis propuesta,
\[
p\approx2,\qquad L\approx1.2,\qquad
|Q(1/16)-L|\approx0.000625<0.001.
\]
El patrón apoya esa estimación condicional, pero no garantiza el error respecto de \(y(T)\). Las dos salidas \(1.200\) sólo coinciden a la precisión mostrada; no prueban que los errores sean iguales o nulos. Más cifras permitirían estudiar el refinamiento; una referencia fiable o una cota de error justificada permitirían comprobar la tolerancia.
